% ─────────────────────────────────────────────────────────────────────────────
%  Cold Email Workflow — Resume Template
%  Author placeholder: fill {{STUDENT_NAME}}, {{EMAIL}}, etc. before compiling.
%
%  Injection marker for Agent3:
%    %%PROJECTS_BEGIN%%
\resumeProjectHeading
  {\textbf{Real-Time 3D Object Detection (LiDAR)} $|$
   \emph{PyTorch, CUDA, ROS2, TensorRT}}{2022.09 -- 2023.06}
  \resumeItemListStart
    \resumeItem{Engineered a lightweight PointPillars-based 3D object detection model for autonomous driving, processing 32-beam LiDAR point clouds in real-time at \textbf{25\,FPS} on an NVIDIA Jetson AGX Xavier.}
    \resumeItem{Designed a multi-scale voxel feature pyramid (MVFP) to enhance 3D perception, improving car AP by \textbf{3.2\%} on the challenging KITTI \textit{hard} test split for autonomous vehicles.}
    \resumeItem{Optimized deep learning inference by developing a custom CUDA kernel for pillar scatter, reducing GPU memory-bandwidth usage by \textbf{28\%} and achieving a 1.4$\times$ throughput improvement for real-time processing.}
    \resumeItem{Integrated the 3D object detector into a ROS2 perception stack for autonomous vehicles, validating end-to-end performance on a real vehicle with zero missed detections over a 500\,m test route.}
  \resumeItemListEnd

\resumeProjectHeading
  {\textbf{Multilingual Low-Resource NER} $|$
   \emph{PyTorch, HuggingFace Transformers, W\&B}}{2023.09 -- 2024.05}
  \resumeItemListStart
    \resumeItem{Designed a deep learning pipeline for Named Entity Recognition (NER), leveraging cross-lingual transfer learning with mBERT and XLM-R backbone models.}
    \resumeItem{Developed a novel data-augmentation strategy, improving macro-F1 performance by \textbf{4.7\%} over the supervised baseline for deep learning models.}
    \resumeItem{Conducted systematic ablation studies in a distributed training setup (4$\times$ A100 GPUs), optimizing deep learning model training wall-clock time by 38\% via gradient checkpointing.}
    \resumeItem{Open-sourced robust training code and evaluation scripts, demonstrating strong software engineering practices and reproducibility in deep learning research.}
  \resumeItemListEnd

\resumeProjectHeading
  {\textbf{Fault-Tolerant Distributed Key-Value Store} $|$
   \emph{Go, gRPC, Raft Consensus, Kubernetes}}{2023.02 -- 2023.08}
  \resumeItemListStart
    \resumeItem{Engineered a linearizable distributed key-value store from scratch, implementing the Raft consensus protocol for robust leader election and log replication.}
    \resumeItem{Achieved \textbf{99.95\% availability} under simulated Byzantine faults and network partitions, sustaining \textbf{12,000 ops/s} at p99 latency $<$ 8\,ms under peak load.}
    \resumeItem{Designed an efficient Bloom-filter and compaction pipeline, reducing on-disk storage overhead by \textbf{40\%} for write-heavy workloads.}
    \resumeItem{Deployed the system on Kubernetes and instrumented with Prometheus/Grafana for real-time visibility into critical operational metrics, ensuring system reliability.}
  \resumeItemListEnd
%%PROJECTS_END%%
% ─────────────────────────────────────────────────────────────────────────────

\documentclass[letterpaper,11pt]{article}

\usepackage{latexsym}
\usepackage[empty]{fullpage}
\usepackage{titlesec}
\usepackage{marvosym}
\usepackage[usenames,dvipsnames]{color}
\usepackage{verbatim}
\usepackage{enumitem}
\usepackage[hidelinks]{hyperref}
\usepackage{fancyhdr}
\usepackage[english]{babel}
\usepackage{tabularx}
\usepackage{fontawesome5}
\input{glyphtounicode}

\pagestyle{fancy}
\fancyhf{}
\fancyfoot{}
\renewcommand{\headrulewidth}{0pt}
\renewcommand{\footrulewidth}{0pt}

\addtolength{\oddsidemargin}{-0.5in}
\addtolength{\evensidemargin}{-0.5in}
\addtolength{\textwidth}{1in}
\addtolength{\topmargin}{-.5in}
\addtolength{\textheight}{1.0in}

\urlstyle{same}
\raggedbottom
\raggedright
\setlength{\tabcolsep}{0in}

\titleformat{\section}{
  \vspace{-4pt}\scshape\raggedright\large
}{}{0em}{}[\color{black}\titlerule \vspace{-5pt}]

\pdfgentounicode=1

% ─── Custom Commands ──────────────────────────────────────────────────────────

\newcommand{\resumeItem}[1]{
  \item\small{#1 \vspace{-2pt}}
}

\newcommand{\resumeSubheading}[4]{
  \vspace{-2pt}\item
    \begin{tabular*}{0.97\textwidth}[t]{l@{\extracolsep{\fill}}r}
      \textbf{#1} & #2 \\
      \textit{\small#3} & \textit{\small #4} \\
    \end{tabular*}\vspace{-7pt}
}

\newcommand{\resumeProjectHeading}[2]{
    \item
    \begin{tabular*}{0.97\textwidth}{l@{\extracolsep{\fill}}r}
      \small#1 & #2 \\
    \end{tabular*}\vspace{-7pt}
}

\newcommand{\resumeItemListStart}{\begin{itemize}}
\newcommand{\resumeItemListEnd}{\end{itemize}\vspace{-5pt}}
\newcommand{\resumeSubHeadingListStart}{\begin{itemize}[leftmargin=0.15in, label={}]}
\newcommand{\resumeSubHeadingListEnd}{\end{itemize}}

% =============================================================================
\begin{document}

% ─── Header ───────────────────────────────────────────────────────────────────
\begin{center}
  \textbf{\Huge \scshape {{STUDENT_NAME}}} \\ \vspace{1pt}
  \small
  \faPhone\ {{PHONE}} $|$
  \href{mailto:{{EMAIL}}}{\faEnvelope\ \underline{{{EMAIL}}}} $|$
  \href{{{GITHUB_URL}}}{\faGithub\ \underline{{{GITHUB_HANDLE}}}} $|$
  \href{{{LINKEDIN_URL}}}{\faLinkedin\ \underline{{{LINKEDIN_HANDLE}}}}
\end{center}

% ─── Education ────────────────────────────────────────────────────────────────
\section{Education}
\resumeSubHeadingListStart
  \resumeSubheading
    {{{UNIVERSITY}}}{{{LOCATION}}}
    {{{DEGREE}}}{{{START_DATE}} -- {{END_DATE}}}
    \resumeItemListStart
      \resumeItem{GPA: {{GPA}} / 4.0 \quad
                  Relevant Courses: {{COURSES}}}
    \resumeItemListEnd
\resumeSubHeadingListEnd

% ─── Research Experience ──────────────────────────────────────────────────────
\section{Research Experience}
\resumeSubHeadingListStart
  \resumeSubheading
    {{{RESEARCH_LAB}}}{{{RESEARCH_LOCATION}}}
    {Research Assistant}{{{RESEARCH_START}} -- {{RESEARCH_END}}}
    \resumeItemListStart
      \resumeItem{{{RESEARCH_DESC_1}}}
      \resumeItem{{{RESEARCH_DESC_2}}}
    \resumeItemListEnd
\resumeSubHeadingListEnd

% ─── Projects ─────────────────────────────────────────────────────────────────
\section{Projects}
\resumeSubHeadingListStart
%%PROJECTS_BEGIN%%
\resumeProjectHeading
  {\textbf{Real-Time 3D Object Detection (LiDAR)} $|$
   \emph{PyTorch, CUDA, ROS2, TensorRT}}{2022.09 -- 2023.06}
  \resumeItemListStart
    \resumeItem{Engineered a lightweight PointPillars-based 3D object detection model for autonomous driving, processing 32-beam LiDAR point clouds in real-time at \textbf{25\,FPS} on an NVIDIA Jetson AGX Xavier.}
    \resumeItem{Designed a multi-scale voxel feature pyramid (MVFP) to enhance 3D perception, improving car AP by \textbf{3.2\%} on the challenging KITTI \textit{hard} test split for autonomous vehicles.}
    \resumeItem{Optimized deep learning inference by developing a custom CUDA kernel for pillar scatter, reducing GPU memory-bandwidth usage by \textbf{28\%} and achieving a 1.4$\times$ throughput improvement for real-time processing.}
    \resumeItem{Integrated the 3D object detector into a ROS2 perception stack for autonomous vehicles, validating end-to-end performance on a real vehicle with zero missed detections over a 500\,m test route.}
  \resumeItemListEnd

\resumeProjectHeading
  {\textbf{Multilingual Low-Resource NER} $|$
   \emph{PyTorch, HuggingFace Transformers, W\&B}}{2023.09 -- 2024.05}
  \resumeItemListStart
    \resumeItem{Designed a deep learning pipeline for Named Entity Recognition (NER), leveraging cross-lingual transfer learning with mBERT and XLM-R backbone models.}
    \resumeItem{Developed a novel data-augmentation strategy, improving macro-F1 performance by \textbf{4.7\%} over the supervised baseline for deep learning models.}
    \resumeItem{Conducted systematic ablation studies in a distributed training setup (4$\times$ A100 GPUs), optimizing deep learning model training wall-clock time by 38\% via gradient checkpointing.}
    \resumeItem{Open-sourced robust training code and evaluation scripts, demonstrating strong software engineering practices and reproducibility in deep learning research.}
  \resumeItemListEnd

\resumeProjectHeading
  {\textbf{Fault-Tolerant Distributed Key-Value Store} $|$
   \emph{Go, gRPC, Raft Consensus, Kubernetes}}{2023.02 -- 2023.08}
  \resumeItemListStart
    \resumeItem{Engineered a linearizable distributed key-value store from scratch, implementing the Raft consensus protocol for robust leader election and log replication.}
    \resumeItem{Achieved \textbf{99.95\% availability} under simulated Byzantine faults and network partitions, sustaining \textbf{12,000 ops/s} at p99 latency $<$ 8\,ms under peak load.}
    \resumeItem{Designed an efficient Bloom-filter and compaction pipeline, reducing on-disk storage overhead by \textbf{40\%} for write-heavy workloads.}
    \resumeItem{Deployed the system on Kubernetes and instrumented with Prometheus/Grafana for real-time visibility into critical operational metrics, ensuring system reliability.}
  \resumeItemListEnd
%%PROJECTS_END%%
\resumeSubHeadingListEnd

% ─── Technical Skills ─────────────────────────────────────────────────────────
\section{Technical Skills}
\begin{itemize}[leftmargin=0.15in, label={}]
  \small{
    \item \textbf{Languages:}  {{SKILLS_LANGUAGES}}  \\
    \item \textbf{Frameworks:} {{SKILLS_FRAMEWORKS}} \\
    \item \textbf{Tools:}      {{SKILLS_TOOLS}}
  }
\end{itemize}

% ─── Publications (optional) ─────────────────────────────────────────────────
% \section{Publications}
% \begin{itemize}[leftmargin=0.15in, label={}]
%   \small{\item{Your publication here.}}
% \end{itemize}

\end{document}
